% ===============================================================
% Section III  --  Methodology
% Owner: Model Development Lead (integrated)
% ===============================================================
\section{Methodology}
\label{sec:methodology}

This project adapts Gemma~3~4B, a pretrained small language model running
locally via Ollama, for Shakespeare question answering using
retrieval-augmented generation (RAG) as the primary adaptation strategy.
Rather than modifying the model's weights, the system conditions
responses on passages retrieved from the Shakespeare corpus at query
time. Three systems were implemented and compared: a prompt-only baseline
that relies solely on the model's pretrained knowledge; a standard RAG
system that retrieves relevant utterances from an indexed corpus before
generating; and an enhanced RAG system that extends retrieval with
cross-encoder reranking and scene-level diversity filtering. All three
share the same language model and hardware, isolating retrieval and
reranking as the variables under study.

\subsection{Baseline System}
The baseline performs no retrieval. A user question is formatted with a
prompt template and sent directly to Gemma~3~4B, which answers solely
from its pretrained weights with no access to the corpus, retrieved
passages, or metadata. This is a prompt-only baseline rather than a
retrieval-only one, a deliberate choice: a no-retrieval baseline produces
a clean separation between the two systems, so that the contribution of
retrieval is directly measurable. When the model lacks retrieved context,
failures such as hallucination, vagueness, and unsupported claims become
visible and attributable to the \emph{absence} of retrieval rather than
to poor retrieval quality. Because the baseline uses the same model,
hardware, and prompt structure as the improved system, any difference in
output quality is attributable to retrieval, not to model capability.

\subsection{RAG-Based System}

\subsubsection{Standard RAG}
The improved system retrieves from the
\texttt{shakespeare\_utterances} ChromaDB collection built by the data
pipeline. Each utterance was indexed using an enriched embedding string
combining speaker, utterance text, scene summary, and keywords, giving
the embedding model enough semantic context to represent short archaic
lines meaningfully.

At query time the user's question is embedded with
\texttt{sentence-transformers/all-MiniLM-L6-v2}, which maps text into a
384-dimensional vector space, and compared against the indexed utterance
vectors via ChromaDB's approximate nearest-neighbour search. Because the
embedding model was trained on large modern-English corpora, it captures
semantic proximity between contemporary query language and archaic
phrasing: a query such as \emph{``Why does Macbeth kill Duncan?''} can
surface relevant passages even when the words \emph{kill} or \emph{Duncan}
do not appear in them, because terms like \emph{kill}, \emph{murder}, and
\emph{do the deed} occupy nearby regions of the space regardless of how
they are written. The enriched embedding strings, with modern-English
scene summaries and keywords alongside the archaic text, further bridge
this vocabulary gap. The top-$k$ results are returned with full metadata
(play, act, scene, speaker, scene summary), and retrieved passages are
always displayed to the user before the generated answer, in every mode,
guaranteeing transparency about the context the model was given.

The system supports four interaction modes. In \texttt{qa} and
\texttt{concept} modes, retrieved utterances are formatted into a
structured context block, with provenance labels and scene context, and
injected into a prompt template alongside the question. In
\texttt{evidence} mode, generation is skipped and only retrieved passages
are shown. In \texttt{stylised} mode, a separate template requests a
Shakespearean-style response limited to 150 words, explicitly labelled as
creative rather than factual. Prompt templates are external text files,
decoupling prompt design from pipeline logic.

\subsubsection{Enhanced RAG (Reranking)}
The enhanced pipeline adds a two-stage retrieve-then-rerank step. In the
first stage, ChromaDB retrieves a broad set of candidates (15 by default)
using the bi-encoder, where query and document vectors are computed
independently and compared by cosine similarity. This is fast but can
miss fine-grained relevance because the two texts are never compared
directly. In the second stage, a cross-encoder
(\texttt{cross-encoder/ms-marco-MiniLM-L-6-v2}) jointly processes each
query--document pair, attending to interactions between query and
document terms, and produces more accurate relevance scores at higher
computational cost, which is why it is applied only to the shortlisted
candidates. Candidates are re-sorted by cross-encoder score before the
top-5 are selected for generation.

A scene-level diversity constraint of at most two results per scene is
applied after reranking. This prevents a single highly relevant scene
from dominating the context, which matters for multi-act questions such
as Hamlet's delayed revenge, where evidence spans several scenes. The
constraint is a known limitation: for focused single-scene questions it
may discard highly relevant passages in favour of weaker ones elsewhere.
A minimum-coverage rule, guaranteeing at least one result per act rather
than capping per scene, is identified as a direction for improvement.

\subsection{Model Choice and Justification}
Gemma~3~4B was selected as the generation model for both the baseline and
the improved systems, running locally via Ollama.
Table~\ref{tab:model_comparison} summarises the candidate models
considered.

\begin{table}[ht]
\centering
\caption{Candidate model comparison. RAM figures from official Ollama
model pages~\cite{ollama_gemma3,localaimaster_ram}. ``RAG eval genuine''
reflects whether the model's limited domain knowledge makes retrieval
genuinely necessary for the comparison.}
\label{tab:model_comparison}
\begin{tabular}{llcc}
\toprule
\textbf{Model} & \textbf{Type} & \textbf{RAM} & \textbf{RAG eval} \\
\midrule
\textbf{Gemma 3 4B} & Local & $\sim$2.5\,GB & 5/5 \\
Mistral 7B Q4       & Local & $\sim$4.3\,GB & 5/5 \\
Gemma 4 E4B         & Local & $\sim$8\,GB    & 5/5 \\
Phi-4 Mini          & Local & $\sim$2.3\,GB & 5/5 \\
TinyLlama 1.1B      & Local & $\sim$0.7\,GB & 5/5 \\
Gemini Flash        & API   & 0\,GB         & 1/5 \\
\bottomrule
\end{tabular}
\end{table}

The deciding hardware constraint was the group's weakest machine at
8\,GB RAM. At roughly 2.5\,GB under 4-bit
quantisation~\cite{google_gemma3_qat}, Gemma~3~4B runs comfortably on
every member's laptop without GPU resources. Mistral~7B~Q4 was eliminated
despite comparable quality because its 4.3\,GB footprint leaves
insufficient headroom on 8\,GB systems~\cite{mistral_vram}; Gemma~4~E4B
sits at the 8\,GB ceiling~\cite{google_gemma4}; Phi-4~Mini fits but is
optimised for reasoning and STEM rather than stylistic
generation~\cite{phi4_mini_benchmark}; and TinyLlama~1.1B struggles with
nuanced synthesis~\cite{tinyllama_eval}.

The choice of a local model over a hosted API was deliberate. Gemma~3~4B
has limited Shakespeare knowledge, so retrieval quality directly and
visibly affects answer quality, which makes the baseline-versus-RAG
comparison meaningful. A powerful hosted model such as Gemini Flash
already answers many Shakespeare questions correctly without context,
which would make the comparison inconclusive and undermine the purpose of
the RAG pipeline.

Fine-tuning was considered and set aside. Adapter methods such as LoRA
shift the output distribution toward training patterns rather than
injecting new factual knowledge, so fine-tuning on Shakespeare would
improve stylistic generation but not factual reliability on plot
questions; it would also blur the experimental signal, since a fine-tuned
baseline would differ from the RAG system for reasons unrelated to
retrieval. Prompt engineering and retrieval are therefore the sole
adaptation strategies. Gemma~3~4B also follows instructions reliably and
produces language of sufficient quality for both beginner-friendly
explanation and stylised generation, making it appropriate for all four
interaction modes.
